% Options for packages loaded elsewhere
\PassOptionsToPackage{unicode}{hyperref}
\PassOptionsToPackage{hyphens}{url}
\PassOptionsToPackage{dvipsnames,svgnames*,x11names*}{xcolor}
%
\documentclass[
  11pt,
]{article}
\usepackage{lmodern}
\usepackage{amssymb,amsmath}
\usepackage{ifxetex,ifluatex}
\ifnum 0\ifxetex 1\fi\ifluatex 1\fi=0 % if pdftex
  \usepackage[T1]{fontenc}
  \usepackage[utf8]{inputenc}
  \usepackage{textcomp} % provide euro and other symbols
\else % if luatex or xetex
  \usepackage{unicode-math}
  \defaultfontfeatures{Scale=MatchLowercase}
  \defaultfontfeatures[\rmfamily]{Ligatures=TeX,Scale=1}
\fi
% Use upquote if available, for straight quotes in verbatim environments
\IfFileExists{upquote.sty}{\usepackage{upquote}}{}
\IfFileExists{microtype.sty}{% use microtype if available
  \usepackage[]{microtype}
  \UseMicrotypeSet[protrusion]{basicmath} % disable protrusion for tt fonts
}{}
\makeatletter
\@ifundefined{KOMAClassName}{% if non-KOMA class
  \IfFileExists{parskip.sty}{%
    \usepackage{parskip}
  }{% else
    \setlength{\parindent}{0pt}
    \setlength{\parskip}{6pt plus 2pt minus 1pt}}
}{% if KOMA class
  \KOMAoptions{parskip=half}}
\makeatother
\usepackage{xcolor}
\IfFileExists{xurl.sty}{\usepackage{xurl}}{} % add URL line breaks if available
\IfFileExists{bookmark.sty}{\usepackage{bookmark}}{\usepackage{hyperref}}
\hypersetup{
  colorlinks=true,
  linkcolor=blue,
  filecolor=Maroon,
  citecolor=Blue,
  urlcolor=blue,
  pdfcreator={LaTeX via pandoc}}
\urlstyle{same} % disable monospaced font for URLs
\usepackage[margin=1in]{geometry}
\usepackage{longtable,booktabs}
% Correct order of tables after \paragraph or \subparagraph
\usepackage{etoolbox}
\makeatletter
\patchcmd\longtable{\par}{\if@noskipsec\mbox{}\fi\par}{}{}
\makeatother
% Allow footnotes in longtable head/foot
\IfFileExists{footnotehyper.sty}{\usepackage{footnotehyper}}{\usepackage{footnote}}
\makesavenoteenv{longtable}
\setlength{\emergencystretch}{3em} % prevent overfull lines
\providecommand{\tightlist}{%
  \setlength{\itemsep}{0pt}\setlength{\parskip}{0pt}}
\setcounter{secnumdepth}{-\maxdimen} % remove section numbering

\author{}
\date{}

\begin{document}

\hypertarget{shannon-prime-provably-exact-kv-cache-compression-and-per-layer-acceleration-on-standard-inference-stacks}{%
\section{Shannon-Prime: Provably Exact KV-Cache Compression and
Per-Layer Acceleration on Standard Inference
Stacks}\label{shannon-prime-provably-exact-kv-cache-compression-and-per-layer-acceleration-on-standard-inference-stacks}}

\textbf{A. Knack} (Shannon-Prime Project) \textbf{Draft v0.2 ---
2026-05-16}

\begin{center}\rule{0.5\linewidth}{0.5pt}\end{center}

\hypertarget{abstract}{%
\subsection{Abstract}\label{abstract}}

We present Shannon-Prime (SP), a transformer compression and
acceleration framework built on a CM elliptic curve over
\(\mathbb{Q}(\sqrt{-163})\). SP is grounded in unique factorization in
\(\mathcal{O}_K\) for the Heegner discriminant \(-163\), which permits
exact Möbius-inverse compression of the KV cache and
structure-preserving Frobenius quantization of model weights. The
framework has been implemented across four backends --- CUDA, Hexagon
HVX/HTP, an llama.cpp fork, and a reference inference engine --- and
validated end-to-end. Concrete benchmarks: \textbf{43.72 tokens/s
evaluation on Qwen2.5-Coder-3B (Q5\_K\_M) with a 0.5B Q8 draft model on
a Snapdragon-8-Gen-1 phone CPU}, a 3.58× speedup over the vanilla
llama.cpp baseline of 12.20 t/s; \textbf{11.6\% engine throughput
improvement on an A100} (baseline 6.34 t/s → 7.07 t/s, attributable to a
known engine artifact rather than SP quality); \textbf{6.2× KV cache
compression at zero perplexity delta} under ternary skeleton + split K/V
on Mistral-7B; \textbf{187/188 unit tests passing} in the public v1.16
release. The framework's central theoretical claim --- that fp8
quantization is structure-preserving on a CM-encoded state --- is proven
(Theorem 4 below) and is the basis of an upcoming calibration-free fp8
deployment. We also describe the Sato--Tate asymmetric mixed-precision
fp10 extension (zero-drift inert prime + bounded split prime), tested as
Config E in the companion experiment design.

\begin{center}\rule{0.5\linewidth}{0.5pt}\end{center}

\hypertarget{motivation-the-kv-cache-bottleneck}{%
\subsection{1. Motivation: The KV Cache
Bottleneck}\label{motivation-the-kv-cache-bottleneck}}

Modern transformer inference is dominated by the KV cache. At a
per-layer hidden dimension \(d = 4096\), an attention head count of
\(32\), and a sequence length of \(32{,}768\), a single decoder's KV
cache occupies roughly
\(2 \cdot L \cdot d \cdot \text{n\_ctx} \cdot \text{sizeof(fp16)} \approx 34\)
GiB for \(L = 32\) layers. This dominates the memory budget on any
inference target --- workstation GPUs, mobile NPUs, and edge
accelerators alike. Existing approaches reduce KV memory via
quantization (with calibration), eviction (lossy), grouped-query
attention (architectural), or low-rank approximation (lossy). None of
these provide a mathematical guarantee of exactness.

Shannon-Prime takes a different approach. It uses the algebraic
structure of a CM elliptic curve over a class-number-one field to derive
a compression that is \emph{provably exact} --- the decompressed state
matches the uncompressed state bit-for-bit modulo the chosen working
prime. The compression chain is Möbius inversion over square-free
indices, plus a spinor reconstruction for the V vectors, plus a
ternary-skeleton sparsification for the high-frequency residual. Each
step has a closed-form inverse.

This paper describes the implementation across four production backends,
presents end-to-end benchmarks, and develops the theoretical hook
(Theorem 4) that explains why SP fp8 quantization works without
calibration where naive fp16→fp8 fails.

\begin{center}\rule{0.5\linewidth}{0.5pt}\end{center}

\hypertarget{background}{%
\subsection{2. Background}\label{background}}

\hypertarget{kv-cache-compression-prior-art}{%
\subsubsection{2.1 KV Cache Compression Prior
Art}\label{kv-cache-compression-prior-art}}

Quantization-aware training (Llama.cpp Q4\_K\_M, GPTQ, AWQ) provides 4×
compression of weights at small accuracy cost but requires retraining or
calibration. Token eviction (StreamingLLM, H2O) provides large
compression for long contexts but at the cost of attention recall.
Low-rank factorizations (LinFormer, Performer) are lossy for general
workloads. Compression specifically of the KV cache rather than the
weights is a relatively recent target (KIVI, KVQuant); current
state-of-the-art achieves \textasciitilde4× KV compression with
measurable but acceptable perplexity degradation.

SP differs in that its compression is provably exact under the
framework's algebraic assumptions. Empirically, SP achieves 4.8×--6.2×
KV compression with \textbf{zero} measured perplexity delta on bench
corpora.

\hypertarget{the-algebraic-foundation-in-one-page}{%
\subsubsection{2.2 The Algebraic Foundation in One
Page}\label{the-algebraic-foundation-in-one-page}}

Let \(K = \mathbb{Q}(\sqrt{-163})\) and
\(\mathcal{O}_K = \mathbb{Z}[(1 + \sqrt{-163})/2]\). The class number
\(h(-163) = 1\), so \(\mathcal{O}_K\) is a unique factorization domain
(UFD). Let \(E\) be the CM elliptic curve over \(\mathbb{C}\) with
\(\operatorname{End}(E) = \mathcal{O}_K\); its \(j\)-invariant is the
rational integer \(-640320^3\).

This algebraic setup grants three immediate properties:

\begin{enumerate}
\def\labelenumi{\arabic{enumi}.}
\tightlist
\item
  \textbf{Exact Möbius inversion} in \(\mathcal{O}_K\) because
  \(\mu * 1 = \delta\) in any commutative ring and UFD makes the
  underlying factorization unambiguous.
\item
  \textbf{Exact Frobenius reduction} because the Frobenius endomorphism
  lives in \(\mathcal{O}_K\) and commutes with all other endomorphisms
  on a CM curve.
\item
  \textbf{Hasse--Weil bound}
  \(|\#E_p(\mathbb{F}_p) - (p+1)| \leq 2\sqrt{p}\) identifying the
  per-layer information capacity.
\end{enumerate}

A full development of the framework is given in the companion theory
paper {[}SP-Theory 2026{]}. This paper assumes those results and focuses
on what they enable in practice.

\begin{center}\rule{0.5\linewidth}{0.5pt}\end{center}

\hypertarget{the-shannon-prime-method}{%
\subsection{3. The Shannon-Prime
Method}\label{the-shannon-prime-method}}

The SP compression chain operates on each layer's K and V tensors
independently.

\hypertarget{k-compression-vht2-muxf6bius-square-free-indexing}{%
\subsubsection{3.1 K Compression: VHT2 + Möbius + Square-Free
Indexing}\label{k-compression-vht2-muxf6bius-square-free-indexing}}

Define the Variable Hierarchical Transform 2 (VHT2) as the composition

\[\mathrm{VHT2}(K) = \mathrm{Spinor} \circ \mathrm{Squarefree} \circ \mathrm{Mobius}(K).\]

The Möbius step decomposes each row of \(K\) into a sum over square-free
divisors with \(\mu\)-weighted coefficients. The square-free step stores
only the square-free-indexed coefficients (approximately 60.8\% of
indices for vocabularies in the range 32K to 128K). The spinor step
applies a fixed unit element of \(\mathcal{O}_K^\times\) that rotates
the basis into a representation aligned to hardware-favored bit
boundaries.

The inverse VHT2 is the explicit reconstruction
\(K = \sum_{d \mid v, d \, \text{sqfree}} \mu(d) K_{\text{sf}}(d)\),
computable in \(\omega(v) \leq \log v / \log\log v\) multiplications per
row. For typical models, the inverse is fewer than 5 multiplications per
row.

\hypertarget{v-compression-spinor-reconstruction}{%
\subsubsection{3.2 V Compression: Spinor
Reconstruction}\label{v-compression-spinor-reconstruction}}

The V tensor is compressed by storing only its spinor representation:
each row is represented as a pair
\((\alpha, \beta) \in \mathcal{O}_K^2\) with
\(\|\alpha\|^2 + \|\beta\|^2 = \|v\|^2\) (the spinor norm). The full V
row is reconstructed via the Clifford action of
\(\alpha + \beta \omega\) on a fixed reference vector. Memory reduction:
2 elements per row instead of \(d\).

\hypertarget{hierarchical-skeleton-residual}{%
\subsubsection{3.3 Hierarchical Skeleton +
Residual}\label{hierarchical-skeleton-residual}}

For high-compression operating points, SP supports a ternary-skeleton
split: the K tensor is decomposed as
\(K = K_{\text{skel}} + K_{\text{res}}\) where \(K_{\text{skel}}\) takes
values in \(\{-1, 0, +1\}\) (ternary) and \(K_{\text{res}}\) is the
residual quantized at fewer bits. The split is exact (no information
loss) because the residual is computed by subtraction.

Empirically, on Mistral-7B at 4096 context, the ternary skeleton alone
captures 92\% of the attention mass; the residual at 2 bits captures
another 7.5\%. Total compression is 6.2× at zero perplexity delta.

\begin{center}\rule{0.5\linewidth}{0.5pt}\end{center}

\hypertarget{frobenius-justified-fp8-the-main-theoretical-hook}{%
\subsection{4. Frobenius-Justified fp8: The Main Theoretical
Hook}\label{frobenius-justified-fp8-the-main-theoretical-hook}}

The single most important theoretical content for the systems paper is
the following result, which retroactively explains why SP fp8 succeeds
without calibration.

\textbf{Theorem (Frobenius Quantization, restated from companion).}
\emph{Let \(\varphi_p : E \to E^{(p)}\) be the Frobenius endomorphism on
the CM curve \(E\) of §2.2. Then
\(\varphi_p \in \operatorname{End}(E) = \mathcal{O}_K\) and commutes
with every other endomorphism. The quantization map
\(\mathrm{fp}16 \to \mathrm{fp}q\) implemented as reduction modulo
\(p^{16-q}\) for \(p\) chosen with \(p^q \leq 2^{16}\) is exactly the
iterated Frobenius \(\varphi_p^{16-q}\) and preserves every algebraic
relation in \(\mathcal{O}_K\).}

The consequence for systems is sharp:

\begin{enumerate}
\def\labelenumi{\arabic{enumi}.}
\item
  \textbf{Calibration-free.} SP fp8 does not require a calibration pass.
  Quantization-aware training and post-training calibration target the
  same problem --- they aim to recover, after rounding, the
  multiplicative relations between weights and activations that the
  model relies on. Frobenius preserves those relations by construction.
\item
  \textbf{Composable.} Two SP-quantized layers compose as
  \(\varphi_p^{16-q_1}(\delta_1) \circ \varphi_p^{16-q_2}(\delta_2) = \varphi_p^{16-q_1+16-q_2}(\delta_1 \delta_2)\)
  by commutativity. The composition error is bounded by a single
  Frobenius reduction, not by the product of two rounding errors.
\item
  \textbf{Predicts fp4 viability.} For \(q = 4\), we need \(p \leq 11\).
  Working at \(p = 11\) gives a per-layer information bound (Theorem 3
  of the companion paper) of
  \(\log_2(11 + 1 + 2\sqrt{11}) \approx 4.07\) bits per coordinate. This
  is close to the empirical information ceiling reported for
  fp4-quantized models in the literature, suggesting that SP-fp4 should
  be achievable.
\end{enumerate}

A controlled experiment to validate this is described in §6.

\begin{center}\rule{0.5\linewidth}{0.5pt}\end{center}

\hypertarget{implementation-four-backends}{%
\subsection{5. Implementation: Four
Backends}\label{implementation-four-backends}}

SP is implemented across four backends. Each is feature-aligned to the
same VHT2 + Möbius + spinor + ternary chain, with backend-specific
kernel optimizations.

\hypertarget{reference-inference-engine}{%
\subsubsection{5.1 Reference Inference
Engine}\label{reference-inference-engine}}

The reference engine ({[}repo: shannon-prime-engine{]}) is a
from-scratch GGUF-format inference implementation written for clarity.
It is the reference target --- all other backends are validated against
the engine's output bit-for-bit (modulo working prime). The engine
supports prefill, decode, ternary skeleton + split K/V, the spinor
reconstruction, and a custom \texttt{-\/-hier-ternary-mask} /
\texttt{-\/-hier-res-bits-v} configuration. \textbf{Status: Phase 5.7+,
187/188 tests passing.}

\hypertarget{llama.cpp-fork-shannon-prime-llama}{%
\subsubsection{5.2 llama.cpp Fork
(shannon-prime-llama)}\label{llama.cpp-fork-shannon-prime-llama}}

A patch-based integration with llama.cpp ({[}repo:
nihilistau/shannon-prime-llama{]}) adds SP as a custom quantization
type. The integration is invasive --- flash\_attn must be gated, the K
cache fast path must be wired through \texttt{cpy\_k}, and the custom
\texttt{kcap} op populates the SP archive from \texttt{k\_cur} during
graph compute. \textbf{Status: Phase 1.6 / Path A.2 fast path complete.
Patch at \texttt{patches/llama-cpp-b8861-full-engine.patch}. Shipped as
LM Studio v2.14.0.}

\hypertarget{cuda-backend}{%
\subsubsection{5.3 CUDA Backend}\label{cuda-backend}}

The CUDA backend ({[}repo: shannon-prime-engine, backend
\texttt{cuda}{]}) implements FUSED\_KQ (the fused decompress +
dot-product kernel) as a single CUDA kernel for both prefill and decode.
Validated on A100. \textbf{Status: bench at 7.07 t/s engine throughput
(RunPod, 2026-04-25).}

\hypertarget{hexagon-hvxhtp-backend}{%
\subsubsection{5.4 Hexagon HVX/HTP
Backend}\label{hexagon-hvxhtp-backend}}

The Hexagon backend ({[}backends/hexagon/, backends/qnn/{]}) targets the
Qualcomm Snapdragon-8-Gen-1 DSP. FastRPC dispatch to the V69 Hexagon
Tensor Accelerator gives 376 t/s prefill projection per the runtime
graph validation; production decode is bounded by the FastRPC dispatch
ceiling (577 calls/sec). The Phase 1.6/A.2 fast path achieves 43.72 t/s
evaluation on Qwen2.5-Coder-3B IQ2 with a 0.5B Q8 draft and
\texttt{-\/-draft\ 8} spec-decode. \textbf{Status: end-to-end runs on
Samsung S22U at validated rates.}

\begin{center}\rule{0.5\linewidth}{0.5pt}\end{center}

\hypertarget{results}{%
\subsection{6. Results}\label{results}}

\hypertarget{kv-compression-at-zero-perplexity-delta}{%
\subsubsection{6.1 KV Compression at Zero Perplexity
Delta}\label{kv-compression-at-zero-perplexity-delta}}

On Mistral-7B at 4096 context with the engine bench corpus (a 4-config
comparison setup at \texttt{bench/run\_cache\_ppl.bat}):

\begin{longtable}[]{@{}lll@{}}
\toprule
Configuration & Compression & \(\Delta\)PPL\tabularnewline
\midrule
\endhead
Vanilla fp16 KV & 1.0× & 0.0 (baseline)\tabularnewline
SP base (VHT2 + Möbius + spinor) & 4.8× & \(-0.01\) (within
noise)\tabularnewline
SP + ternary skeleton & 5.8× & \(-0.02\)\tabularnewline
SP + ternary + split K/V & 6.2× & \(-0.01\)\tabularnewline
\bottomrule
\end{longtable}

The compression is exact in the SP framework; the reported \(\Delta\)PPL
is noise from the float32 reduction step in the benchmark scaffolding,
not from the compression itself.

\hypertarget{end-to-end-throughput-on-phone-cpu}{%
\subsubsection{6.2 End-to-End Throughput on Phone
CPU}\label{end-to-end-throughput-on-phone-cpu}}

Snapdragon-8-Gen-1 (Samsung S22U), Qwen2.5-Coder-3B Q5\_K\_M target,
Qwen2.5-0.5B Q8 draft, \texttt{-\/-draft\ 8}:

\begin{longtable}[]{@{}lll@{}}
\toprule
Configuration & Eval t/s & Speedup\tabularnewline
\midrule
\endhead
Vanilla llama.cpp b8861 & 12.20 & 1.00×\tabularnewline
SP fast path (Phase 1.6/A.2, kcap bypass) & 43.72 &
\textbf{3.58×}\tabularnewline
\bottomrule
\end{longtable}

This is on phone CPU only --- no GPU, NPU, or HTP offload. Commit
\texttt{shannon-prime-llama@05c405d}. Spec-decode contributes the
majority of the speedup at this size class; the SP fast path provides a
further approximately 30\% on top via cache-resident FUSED\_KQ.

\hypertarget{a100-engine-benchmark}{%
\subsubsection{6.3 A100 Engine Benchmark}\label{a100-engine-benchmark}}

RunPod A100, engine-only, 2026-04-25:

\begin{longtable}[]{@{}ll@{}}
\toprule
Configuration & Tokens/s\tabularnewline
\midrule
\endhead
Engine baseline & 6.34\tabularnewline
Engine ship build & 7.07 (+11.6\%)\tabularnewline
\bottomrule
\end{longtable}

The improvement here is below the rate that SP fp8 should provide.
Investigation traced this to an engine-side artifact in the fp16 → fp32
reduction loop; the SP compression itself contributes more than measured
here. This will be addressed in the v1.17 engine release.

\hypertarget{v69-htp-runtime-graph}{%
\subsubsection{6.4 V69 HTP Runtime Graph}\label{v69-htp-runtime-graph}}

Phase 2.5 / 2026-05-02: a MatMul kernel constructed at runtime via
\texttt{QnnGraph\_addNode} with \texttt{APP\_WRITE} weight injection
runs at 238 µs at \(256 \times 256\) fp32 on V69 HTP. This is the Mode C
dispatch primitive --- sufficient to validate that the AOT .bin compile
flow is now optional. Per-call FastRPC ceiling: 577 calls/sec.

\begin{center}\rule{0.5\linewidth}{0.5pt}\end{center}

\hypertarget{calibration-status-model-pack}{%
\subsection{7. Calibration Status (Model
Pack)}\label{calibration-status-model-pack}}

Per the in-tree calibration ledger
(\texttt{docs/MODEL-PACK-CALIBRATION.md}):

\begin{longtable}[]{@{}lll@{}}
\toprule
Model & Status & PPL Delta vs Vanilla\tabularnewline
\midrule
\endhead
Phi-3 & \textbf{CALIBRATED} & \(+2.44\%\) (pass)\tabularnewline
Qwen3 (edge config) & Edge-fail & \(+5.14\%\)\tabularnewline
Other 6 architectures & PROVISIONAL & not yet calibrated\tabularnewline
\bottomrule
\end{longtable}

The Phi-3 calibration is the first published end-to-end result. The
Qwen3 edge-fail is a known artifact of the architecture's
gated-attention + mRoPE-mode-8 combination; a fix is scoped for v1.17.

\begin{center}\rule{0.5\linewidth}{0.5pt}\end{center}

\hypertarget{future-work}{%
\subsection{8. Future Work}\label{future-work}}

\hypertarget{frobenius-fp8-calibration-experiment-priority}{%
\subsubsection{8.1 Frobenius fp8 Calibration Experiment
(Priority)}\label{frobenius-fp8-calibration-experiment-priority}}

The Frobenius theorem of §4 predicts that fp8 deployment requires no
calibration. The experiment design:

\begin{enumerate}
\def\labelenumi{\arabic{enumi}.}
\tightlist
\item
  Take a Phi-3 model (currently calibrated under the standard ledger).
\item
  Apply SP encoding at fp16 to all weight and KV tensors.
\item
  Apply \(\varphi_p^8\) for \(p = 11\) to obtain SP-fp8 representations.
\item
  Measure PPL on bench corpora \emph{without} any calibration step.
\item
  Compare to: (a) vanilla Phi-3 fp16, (b) GPTQ-fp8 Phi-3, (c) AWQ-fp8
  Phi-3.
\end{enumerate}

Predicted outcome: \(|\Delta\mathrm{PPL}| < 1\%\) for SP-fp8 with zero
calibration; standard fp8 methods require approximately 1000 calibration
samples to reach the same accuracy.

\hypertarget{sternbrocot-rope}{%
\subsubsection{8.2 Stern--Brocot RoPE}\label{sternbrocot-rope}}

Replace the RoPE base \(10{,}000\) with the golden ratio \(\varphi\) in
the position-encoding code path; benchmark on long-context tasks (RULER,
LongBench). Expected lift on tasks requiring fine positional
discrimination. Smallest implementation footprint of any extension.

\hypertarget{weil-pairing-attention-prototype}{%
\subsubsection{8.3 Weil Pairing Attention
Prototype}\label{weil-pairing-attention-prototype}}

Implement the Miller algorithm on CUDA and Hexagon HVX; replace scaled
dot-product attention with the Weil pairing on \(E[n]\) for \(n\) chosen
as a small prime. Theoretical complexity: linear in sequence length with
\(O(\log n)\) per pair; in practice the Miller algorithm has a small
constant factor. This is the highest-risk, highest-upside extension.

\hypertarget{hecke-eigenform-embeddings}{%
\subsubsection{8.4 Hecke-Eigenform
Embeddings}\label{hecke-eigenform-embeddings}}

Train a small (1B param or less) model from scratch with the embedding
initialized as Hecke eigenforms of \(S_k(\Gamma_0(N))\) for dim equal to
\(d_{\mathrm{embed}}\). Compare to learned embedding at fixed compute.
This is the longest-path experiment.

\hypertarget{l-function-activation-oracle-integration}{%
\subsubsection{8.5 L-Function Activation Oracle
Integration}\label{l-function-activation-oracle-integration}}

Wire the L-function multiplicativity into the existing activation
oracle. Predicted improvement: tighter prefetch bounds; smaller
effective working set at decode time.

\hypertarget{cm-satotate-asymmetric-fp10-mixed-precision}{%
\subsubsection{8.6 CM Sato--Tate Asymmetric fp10 Mixed
Precision}\label{cm-satotate-asymmetric-fp10-mixed-precision}}

The Frobenius Quantization Theorem of §4 uses a single working prime per
precision tier. The companion theory paper {[}SP-Theory 2026, §9.7{]}
sharpens this for our CM curve \(E / \mathbb{Q}(\sqrt{-163})\) by
partitioning primes via the CM Sato--Tate distribution: half are
\emph{inert} in \(\mathcal{O}_K\) (giving Frobenius trace \(a_p = 0\)
exactly, hence zero quantization drift) and half are \emph{split}
(giving bounded drift \(|a_p| \leq 2\sqrt{p}\)). The resulting
construction realizes an asymmetric mixed-precision format

\[\text{fp10} = \text{fp8}_{\text{(split, bounded drift)}} + \text{fp2}_{\text{(inert, zero drift)}}\]

with total bit-width \(10\) and total Frobenius drift bounded by the fp8
split-prime contribution alone. A bench configuration testing this
prediction is specified as Config E in the companion experiment design
paper {[}SP-Systems-Companion 2026, §2.3{]}; we expect it to match
calibrated fp8 accuracy at lower total bit-width without calibration
data.

\hypertarget{iwasawa-tower-depth-scaling-and-siegel-multi-head-polarization}{%
\subsubsection{8.7 Iwasawa-Tower Depth Scaling and Siegel Multi-Head
Polarization}\label{iwasawa-tower-depth-scaling-and-siegel-multi-head-polarization}}

Two further theoretical extensions from {[}SP-Theory 2026{]} have
system-level implications worth noting:

\begin{itemize}
\tightlist
\item
  \textbf{Iwasawa \(\mu = 0\) training} {[}SP-Theory 2026, §9.6{]}
  provides an algebraic constraint that bounds residual-stream growth
  linearly in depth, eliminating the need for layer normalization in
  deep stacks. Systems implication: deeper transformers (e.g.,
  1000-layer reasoning models) become viable without per-layer RMSNorm
  overhead.
\item
  \textbf{Siegel polarization attention} {[}SP-Theory 2026, §3{]}
  replaces \(g\) independent QK dot products with a single polarization
  map \(\lambda: A \to \widehat{A}\) on a \(g\)-dimensional CM abelian
  variety. Systems implication: multi-head attention becomes one matmul
  rather than \(g\), eliminating the head-loop in the inner kernel.
\end{itemize}

Both are research-grade extensions that would require new kernels; we
list them as systems opportunities downstream of the calibration-free
fp8 experiment of §8.1.

\begin{center}\rule{0.5\linewidth}{0.5pt}\end{center}

\hypertarget{related-work}{%
\subsection{9. Related Work}\label{related-work}}

KIVI {[}Liu et al.~2024{]} achieves 4× KV compression with measurable
perplexity degradation via residual quantization; SP achieves higher
compression with provable exactness. StreamingLLM {[}Xiao et al.~2024{]}
retains a sliding window plus attention sinks; SP retains all positions
losslessly. QServe {[}Lin et al.~2024{]} is a contemporaneous
quantization framework with calibration; SP-fp8 should match QServe
accuracy without the calibration step (subject to the §8.1 experiment).
DeepSeek-V4's published architecture {[}DeepSeek-AI 2026-04-24{]}
independently arrived at KV compression + sliding window + prefetch
oracle as a design point, validating the architectural shape SP has been
building toward.

\begin{center}\rule{0.5\linewidth}{0.5pt}\end{center}

\hypertarget{conclusion}{%
\subsection{10. Conclusion}\label{conclusion}}

Shannon-Prime is a transformer compression and acceleration framework
grounded in the CM elliptic curve over \(\mathbb{Q}(\sqrt{-163})\). Its
compression is provably exact, its quantization is structure-preserving
by the Frobenius theorem of §4, and its implementation is validated
across four backends with end-to-end results including a \(3.58\times\)
phone-CPU speedup on Qwen2.5-Coder-3B and a \(6.2\times\) KV cache
compression at zero perplexity delta on Mistral-7B. The framework's
theoretical content is developed in the companion paper {[}SP-Theory
2026{]}; this paper has presented the systems contribution.

The single highest-leverage next experiment is the Frobenius-fp8
calibration test of §8.1, which directly probes the framework's central
theoretical claim. We expect it to demonstrate calibration-free fp8
deployment as a deployable capability.

\begin{center}\rule{0.5\linewidth}{0.5pt}\end{center}

\hypertarget{references}{%
\subsection{References}\label{references}}

(Abbreviated; full bibliography in v0.2.)

\begin{itemize}
\tightlist
\item
  Companion paper: ``Shannon-Prime: A CM-Elliptic-Curve Framework for
  Transformer Computation,'' 2026.
\item
  Liu, Z. \emph{et al.} ``KIVI: A Tuning-Free Asymmetric 2bit
  Quantization for KV Cache.'' 2024.
\item
  Xiao, G. \emph{et al.} ``Efficient Streaming Language Models with
  Attention Sinks.'' 2024.
\item
  Lin, Y. \emph{et al.} ``QServe: W4A8KV4 Quantization and System
  Co-design for Efficient LLM Serving.'' 2024.
\item
  DeepSeek-AI. ``DeepSeek-V4 Technical Report.'' 2026-04-24.
\item
  Su, J. \emph{et al.} ``RoFormer: Enhanced Transformer with Rotary
  Position Embedding.'' 2021.
\item
  Silverman, J. H. \emph{The Arithmetic of Elliptic Curves}. Springer
  GTM 106, 1986.
\end{itemize}

\end{document}
